\clearpage
\VerdierPhysicalPageAnchor{0138}%
\begingroup
\small
\setlength{\emergencystretch}{2em}%
\setlength{\parskip}{0pt}%
\setlength{\abovedisplayskip}{1pt plus 1pt minus 1pt}%
\setlength{\belowdisplayskip}{1pt plus 1pt minus 1pt}%

\begingroup\itshape
\VerdierSemanticAnchor{ch2-2-2-4-b}
\noindent b) 圏 $\mathcal C$ が加法圏ならば、圏
$\mathcal C(S^{-1})$ は加法圏であり、関手 $Q$ は加法関手である。

\VerdierSemanticAnchor{ch2-2-2-4-c}
\par\noindent c) $\widehat S$ を、$Q$ により $\mathcal C(S^{-1})$ の
同型へ送られる $\mathcal C$ の射の集合とする。
\VerdierSemanticAnchor{ch2-2-2-4-c-saturation}%
集合 $\widehat S$ は $\mathcal C$ の飽和乗法的系である。$S$ を含む
$\mathcal C$ の任意の飽和乗法的系は $\widehat S$ を含む。
\endgroup

\VerdierSemanticAnchor{ch2-2-2-4-proof-b}
\par\noindent 第一の主張は $\mathcal C(S^{-1})$ の普遍性から直ちに従う：
圏 $\bigl(\mathcal C^\circ(S^{-1})\bigr)^\circ$ はこの問題の解である。
第二の主張は定理 (2.2.2) の主張 b)、c) から直ちに従う。
\VerdierSemanticAnchor{ch2-2-2-4-proof-c-characterization}%
$\mathcal C(S^{-1})$ の射の記述から、$\mathcal C$ の射 $f$ が $Q$ に
より同型へ送られるための必要十分条件は、$gf$ と $fh$ が定義されて
$S$ に属するような $\mathcal C$ の二射 $g$、$h$ が存在することだと
直ちに分かる。
\VerdierSemanticAnchor{ch2-2-2-4-proof-c-sm2-sm3}%
したがって、集合 $\widehat S$ は性質 SM2)、SM3) をもつ。
\VerdierSemanticAnchor{ch2-2-2-4-proof-c-sm1-sm4}%
さらに、$\widehat S$ の定義から、この集合は性質 SM1)、SM4) をもつ。
\VerdierSemanticAnchor{ch2-2-2-4-proof-c-conclusion}%
ゆえに $\widehat S$ は $\mathcal C$ の飽和乗法的系である。(2.2.4) の
最後の主張は明らかである。

\VerdierSemanticAnchor{ch2-2-2-5}
\smallskip
\noindent\textbf{2.2.5.} $\mathcal C$ を圏、$S$ を $\mathcal C$ の
乗法的系とする。
\VerdierSemanticAnchor{ch2-2-2-5-localized-category}%
圏 $\mathcal C(S^{-1})$ を、系 $S$ に関する $\mathcal C$ の
\emph{局所化圏}と呼ぶ。
\VerdierSemanticAnchor{ch2-2-2-5-localization-functor}%
関手 $Q\colon\mathcal C\longrightarrow\mathcal C(S^{-1})$ を
\emph{局所化関手}と呼ぶ。
\VerdierSemanticAnchor{ch2-2-2-5-saturation}%
乗法的系 $\widehat S$ を系 $S$ の\emph{飽和化}と呼ぶ。

\VerdierSemanticAnchor{ch2-2-2-6}
\smallskip
\noindent\textbf{定理 2.2.6.}\quad
\begingroup\itshape
\VerdierSemanticAnchor{ch2-2-2-6-setup}%
$\mathcal D$ を三角圏、$S$ を $\mathcal D$ の乗法的系とする (2.1.4)。
$\mathcal D(S^{-1})$ を $\mathcal D$ の局所化圏、
$Q\colon\mathcal D\longrightarrow\mathcal D(S^{-1})$ を局所化関手とする。

\VerdierSemanticAnchor{ch2-2-2-6-a}
\par\noindent a) $\mathcal D(S^{-1})$ 上には、$Q$ が厳密な
$\mathbb Z$-関手となるような厳密な $\mathbb Z$-圏の構造がただ一つ
存在する。

\VerdierSemanticAnchor{ch2-2-2-6-b}
\par\noindent b) $\mathcal D(S^{-1})$ 上には、$Q$ が完全関手となる
ような三角構造がただ一つ存在する。この構造の完全三角は、
$\mathcal D$ の完全三角の $Q$ による像と同型な三角である。

\VerdierSemanticAnchor{ch2-2-2-6-c-factorization}
\par\noindent c) $S$ の射を同型へ送る、三角圏（それぞれアーベル圏）
への任意の完全（それぞれコホモロジー的）関手
$F\colon\mathcal D\longrightarrow\mathcal D'$ は、
$Q\colon\mathcal D\longrightarrow\mathcal D(S^{-1})$ を通じて完全
（それぞれコホモロジー的）関手
$G\colon\mathcal D(S^{-1})\longrightarrow\mathcal D'$ として一意に
分解する。
\VerdierSemanticAnchor{ch2-2-2-6-c-morphism-lifting}%
$F_1,F_2\colon\mathcal D\rightrightarrows\mathcal D'$ を $S$ の射を
同型へ送る二つの完全（それぞれコホモロジー的）関手とし、
$G_1,G_2\colon\mathcal D(S^{-1})\rightrightarrows\mathcal D'$ をそれぞれ
の分解、$m\colon F_1\longrightarrow F_2$ を完全（それぞれ
コホモロジー的）関手の射とする。このとき $m=p\star Q$ を満たす
関手の射 $p\colon G_1\longrightarrow G_2$ がただ一つ存在する。
\endgroup

\VerdierSemanticAnchor{ch2-2-2-6-proof-a}
\par\noindent 第一の主張は定理 (2.2.2) と SM5) から直ちに従う。
\VerdierSemanticAnchor{ch2-2-2-6-proof-c-p139-boundary}%
第三の主張は

\VerdierSemanticAnchor{ch2-2-2-6-p139-boundary}
\endgroup
